\ncertchapter{2}{Mechanical Properties of Fluids}\chlabel{ch:fluids}
% ================================================================

A ship made of steel, a material nearly eight times as dense as water, floats; a steel nail sinks. An
airliner of several hundred tonnes climbs into the sky and stays there. A dam must be built far thicker at
its base than at its top. A straw lets you drink, a syringe delivers medicine, and your heart pushes blood
through thousands of kilometres of vessels. Every one of these is fluid mechanics, the physics of liquids
and gases.

\chref{ch:solids} described a solid by how it resists being deformed: a stress produces a strain, and the
ratio is a modulus. One of those moduli, the shear modulus $G$, measures resistance to a sideways
(shearing) deformation. A \textbf{fluid} is a substance whose shear modulus is zero: at rest it cannot
resist a shearing force at all, however small, so it \emph{flows} and takes the shape of its container.
That single difference reshapes the mechanics. A fluid at rest can only push, and only perpendicular to a
surface, so the one quantity needed to describe its forces is pressure. A fluid in motion carries energy
from place to place, and the work--energy theorem of \chref{ch:energy}, applied to a moving parcel of fluid,
becomes Bernoulli's equation.

The chapter follows that order: fluids at rest (pressure, Pascal's law, buoyancy), then ideal fluids in
motion (continuity and Bernoulli), and finally the two properties that the ideal model leaves out,
viscosity and surface tension. Each model has limits, and they are stated where the model is built: liquids
are treated as incompressible, $g$ as uniform, and flows as steady until we say otherwise.

\begin{motivbox}{Why fluid mechanics matters}
Three ideas do most of the work in this chapter.
\begin{itemize}
\item \textbf{Pressure and Pascal's law build the machines that lift cars.} Because an enclosed liquid
transmits a change of pressure to every part of itself, a small force on a small piston becomes a large
force on a large one: the hydraulic jack, the car brake, the excavator arm and the dentist's chair all use
this one principle. The same law $P=P_0+\rho gh$ tells you why a dam is thick at the base and why divers
feel the water on their ears.
\item \textbf{Buoyancy decides what floats.} Archimedes' principle is why a steel ship floats and a steel
nail sinks, why a hot-air balloon rises, how a submarine chooses its depth, and how a hydrometer measures
the strength of battery acid.
\item \textbf{Bernoulli's equation links speed and pressure.} One energy equation (where a fluid moves
faster, its pressure is lower) accounts for the spray of a perfume atomiser, the reading of a Venturi
flow meter, the speed of water from a hole in a tank, and the pressure difference across an aircraft wing.
\end{itemize}
\end{motivbox}

\section{Pressure in a Fluid}\label{sec:fluid-pressure}

Start with a fluid at rest. Because it cannot resist shear, any force it exerts on a surface must be
\emph{perpendicular} to that surface: a sideways component would be a shear, and the fluid would flow under
the reaction to it. So the force a resting fluid exerts on a small flat area $A$ is fixed by one number, the
force per unit area, whatever the direction the area faces.

\begin{defbox}{Pressure}
The \textbf{pressure} $P$ at a point in a fluid is the magnitude of the normal force per unit area on a
small surface placed at that point:
\[
P=\frac{F_\perp}{A}\qquad\big(\text{unit: }\SI{1}{Pa}=\SI{1}{N/m^2}\big).
\]
Pressure is a scalar: at a point in a fluid at rest it has the same value whichever way the surface faces
(derived below), and the force it produces on any surface points along the surface's normal, into it.
Other units in use: $\SI{1}{atm}=\SI{1.013e5}{Pa}$ (average sea-level air pressure, written $P_0$),
$\SI{1}{bar}=\SI{e5}{Pa}$, and $\SI{1}{mm}$ of mercury, about \SI{133}{Pa}. The pressure measured
relative to the atmosphere, $P-P_0$, is called the \textbf{gauge pressure}; $P$ itself is the
\textbf{absolute pressure}.
\end{defbox}

The claim that pressure has no direction is not obvious, and it is worth proving, because every result
of the chapter leans on it.

\begin{derivbox}{Pressure at a point is the same in every direction}
\stepn{1} \textbf{Isolate a tiny wedge of fluid.} Take a small prism of resting fluid whose cross-section is
a right-angled triangle: a vertical face of area $A_x$, a horizontal face of area $A_y$, and a slanted face
of area $A$ at angle $\theta$ to the horizontal. Geometry gives $A_x=A\sin\theta$ and $A_y=A\cos\theta$.
Let the pressures on the three faces be $P_x$, $P_y$ and $P$; each pushes normally inwards.

\stepn{2} \textbf{Balance horizontally.} The slanted face's push has a horizontal component
$PA\sin\theta$; the vertical face pushes the other way with $P_xA_x=P_xA\sin\theta$. Gravity has no
horizontal part, and the fluid is at rest (Newton's first law, \chref{ch:newton}), so $P_x=P$.

\stepn{3} \textbf{Balance vertically.} The horizontal face pushes up with $P_yA\cos\theta$; the slanted face
pushes down with $PA\cos\theta$; the wedge's weight is $\rho g\times(\text{volume})$, where $\rho$ is the
fluid's density (\si{kg/m^3}) and $g$ the acceleration due to gravity. The areas scale as the square of
the wedge's size but the volume as the cube, so as the wedge shrinks to a point the weight becomes
negligible beside the surface forces, leaving $P_y=P$.

\stepn{4} \textbf{Conclude.} The angle $\theta$ was arbitrary, so $P_x=P_y=P$ for every orientation: at a
point, a fluid at rest pushes equally in all directions.
\end{derivbox}

\inwords{A fluid at rest can only push, never pull or drag sideways, and at any one point it pushes equally
hard in every direction. That is why a diver feels the water squeeze from all sides at once, not just from
above.}

The next question is how the pressure changes from point to point. Within one horizontal level it cannot
change at all (a horizontal difference would push the fluid sideways), but it must grow with depth, because
deeper fluid has more fluid above it to hold up.

\begin{derivbox}{Pressure increases with depth: $P=P_0+\rho gh$}
\stepn{1} \textbf{Choose a column.} In a liquid of uniform density $\rho$ whose free surface is open to the
atmosphere (pressure $P_0$), isolate a vertical cylinder of cross-section $A$ reaching from the surface down
to a depth $h$.

\stepn{2} \textbf{List the forces on it.} The atmosphere pushes down on its top with $P_0A$. Its weight is
$mg=\rho(Ah)g$, downwards. The liquid below pushes up on its base with $PA$, where $P$ is the pressure at
depth $h$. The surrounding liquid also pushes on the curved side, but those pushes are horizontal and, at
each depth, equal from opposite sides, so they cancel.

\stepn{3} \textbf{Balance.} The column is at rest, so the vertical forces sum to zero:
\[
PA=P_0A+\rho Ahg\quad\Longrightarrow\quad\boxed{\;P=P_0+\rho gh\;}
\]

\stepn{4} \textbf{Read it.} Pressure grows linearly with depth. It depends on the vertical depth $h$ only,
not on the width or shape of the container, and it is the same at every point of a horizontal level within
one connected body of liquid. The gauge pressure at depth $h$ is $\rho gh$.
\end{derivbox}

\why{That the pressure depends only on \emph{depth}, not on the shape or amount of liquid, is the
``hydrostatic paradox'': a thin tube of water and a whole lake press equally hard at the same depth. It is
why a dam is built thick at the bottom, where $\rho gh$ is largest; why your ears hurt at the bottom of a
pool; and why a water tower needs to be \emph{tall}, not large, to deliver pressure to a town.}

\begin{trapbox}{the paradox of the base force}
Three vessels of different shapes, with equal base areas $A$, are filled to the same depth $h$. The force
on each base is the same, $\rho ghA$ (plus $P_0A$), even though the vessels hold different weights of
liquid. There is no contradiction: in a vessel that widens upwards, the sloping walls carry part of the
liquid's weight; in one that narrows upwards, the walls push \emph{down} on the liquid and add to the force
on the base. The base force equals the liquid's weight only for straight vertical walls.
\end{trapbox}

The law $P=P_0+\rho gh$ also turns a column of liquid into a pressure gauge, which is how atmospheric
pressure was first measured.

\begin{exbox}{The mercury barometer}\label{ex:fluids-barometer}
A long glass tube, closed at one end, is filled with mercury ($\rho=\SI{13600}{kg/m^3}$) and inverted into
a dish of mercury. The column falls until its top is \SI{760}{mm} above the surface in the dish, leaving a
near-vacuum above it. Find the atmospheric pressure ($g=\SI{9.8}{m/s^2}$).

\stepn{1} The pressure at the top of the column is nearly zero (mercury vapour exerts only about
\SI{0.2}{Pa}). Going down a depth $h=\SI{0.760}{m}$ adds $\rho gh$.

\stepn{2} At the level of the dish's surface the column must have the pressure of the open surface beside
it, which is $P_0$ (same level, same liquid). So
$P_0=\rho gh=13600\times9.8\times0.760\approx\SI{1.013e5}{Pa}$.

\medskip
This is why \SI{1}{atm} is also called \SI{760}{mm} of mercury. A water barometer would need a column
$13.6$ times taller, about \SI{10.3}{m}. \ans{$\approx\SI{1.013e5}{Pa}$}
\end{exbox}

Because pressure grows with depth, the force on a large vertical wall is not simply ``pressure times area'';
the pressure must be added up strip by strip.

\begin{exbox}{The force on a dam}\label{ex:fluids-dam}
A vertical dam wall of width $w=\SI{100}{m}$ holds back water of depth $H=\SI{20}{m}$. Find the total force
of the water on it, beyond what the atmosphere exerts on both faces ($g=\SI{9.8}{m/s^2}$).

\stepn{1} A horizontal strip of the wall at depth $y$, of height $dy$, has area $w\,dy$ and feels gauge
pressure $\rho gy$, so the force on it is $dF=\rho gy\,w\,dy$.

\stepn{2} Add the strips from the surface to the base:
\[
F=\int_0^H\rho g w\,y\,dy=\tfrac12\rho gwH^2=\tfrac12(1000)(9.8)(100)(20)^2\approx\SI{1.96e8}{N}.
\]

\stepn{3} The result is the \emph{average} gauge pressure $\tfrac12\rho gH$ times the area $wH$, because the
pressure rises linearly. Doubling the depth quadruples the force, and the push is concentrated low down:
it acts, on balance, one third of the way up the wall. That is why dams are thickest at the base.
\ans{$\approx\SI{2.0e8}{N}$}
\end{exbox}

\begin{remarkbox}{Where $P=P_0+\rho gh$ stops working}
The derivation assumed a uniform density and a uniform $g$. Both are excellent for liquids over everyday
depths, because a liquid is very hard to compress: at the bottom of the deepest ocean trench, about
\SI{11}{km} down, water is compressed by only about $5\%$. Air is different. Its density is proportional to
its pressure, so it thins with height and the pressure falls not linearly but exponentially, by a factor
$e$ every \SI{8}{km} or so for an atmosphere at a uniform temperature (a first-order decay of the kind in
\chref{ch:toolkit}, Section~\ref{sec:explog}; the gas law it needs, $PV=nRT$, comes in \chref{ch:thermal}, Section~\ref{sec:gaslaw}). Over a
room or a building, though, the change in air pressure is tiny and the linear law is fine.
\end{remarkbox}

\section{Pascal's Law and the Hydraulic Press}\label{sec:pascal}

The depth law contains a further consequence. In $P=P_0+\rho gh$ the pressure at depth $h$ is the surface
pressure \emph{plus} a fixed amount. If the pressure at the surface is raised by $\Delta P$ (say, by
pushing on a piston), then, provided the liquid does not compress, the depth term is unchanged, and the
pressure at every point rises by the same $\Delta P$.

\begin{keybox}{Pascal's law}
A change of pressure applied to an enclosed incompressible fluid is transmitted undiminished to every part
of the fluid and to the walls of its container.
\end{keybox}

The most useful application multiplies force (Figure~\ref{fig:fluids-hydraulic}).

\begin{figure}[htbp]
\centering
\tikzfig{fluids_hydraulic}
\figcaption{fig:fluids-hydraulic}{The hydraulic lift. The pistons sit at the same level of one liquid, so the
pressure under each is the same, $P=F_1/A_1$. On the large piston that pressure gives the much larger
force $F_2=PA_2$. When the small piston moves down $d_1$, the liquid it displaces spreads under the large
piston, which rises only $d_2=d_1A_1/A_2$.}
\end{figure}

\begin{derivbox}{The hydraulic lift multiplies force, not energy}
\stepn{1} A force $F_1$ on the small piston (area $A_1$) creates pressure $F_1/A_1$ beneath it. By Pascal's
law, and because both pistons are at the same level, the same pressure acts under the large piston (area
$A_2$).

\stepn{2} The force on the large piston is therefore
\[
F_2=PA_2=F_1\,\frac{A_2}{A_1},
\]
larger than $F_1$ by the area ratio, which for circular pistons is the square of the ratio of their radii.

\stepn{3} \textbf{Energy check.} The liquid is incompressible, so the volume swept by the small piston,
$A_1d_1$, equals the volume gained under the large one, $A_2d_2$. The work put in is
$F_1d_1=PA_1d_1$; the work delivered is $F_2d_2=PA_2d_2$. These are equal, so the lift gains force at the
cost of distance, exactly as a lever does, and creates no energy.
\end{derivbox}

\begin{exbox}{Lifting a car}\label{ex:fluids-carlift}
A garage lift raises a \SI{1500}{kg} car on a piston of radius \SI{15}{cm}. The small piston has radius
\SI{1.5}{cm}. What force must be applied to the small piston, and how far does the car rise when the small
piston is pushed down \SI{20}{cm}? ($g=\SI{9.8}{m/s^2}$; ignore the weights of the pistons.)

\stepn{1} Area ratio: $A_2/A_1=(15/1.5)^2=100$.

\stepn{2} $F_1=F_2\,A_1/A_2=(1500\times9.8)/100\approx\SI{147}{N}$, about the weight of a \SI{15}{kg} bag.

\stepn{3} $d_2=d_1A_1/A_2=\SI{20}{cm}/100=\SI{2}{mm}$. To raise the car \SI{20}{cm} the small piston must
be pumped through \SI{20}{m} in total, which is why a jack has a valve and many short strokes.
\ans{\SI{147}{N}; \SI{2}{mm}}
\end{exbox}

\begin{trapbox}{Pascal's law needs an incompressible fluid}
If air gets into a car's brake lines, part of the pedal's travel is spent compressing the bubble instead
of moving the brake pistons, and the pedal feels ``spongy''. The law holds for liquids because they barely
compress; in a gas a pressure change is still transmitted, but only after the gas has been squeezed.
\end{trapbox}

\section{Archimedes' Principle and Floating}\label{sec:archimedes}

Why do things feel lighter in water, and why does a steel ship float? Pressure increases with depth, so the
fluid pushes up on the bottom of an immersed body harder than it pushes down on the top. The net upward
push is the \textbf{buoyant force}, $F_B$ (Figure~\ref{fig:fluids-archimedes}).

\begin{figure}[htbp]
\centering
\tikzfig{fluids_archimedes}
\figcaption{fig:fluids-archimedes}{Why a submerged body is pushed up. The pressure pushes are drawn to scale:
the push on the deeper bottom face exceeds the push on the top, while the side pushes at each depth are
equal and opposite and cancel. What remains is the buoyant force $F_B$, which equals the weight of the
fluid the block displaces.}
\end{figure}

\begin{derivbox}{Archimedes' principle}
\stepn{1} \textbf{A block first.} A block of face area $A$ and height $\ell$ is submerged with its top at depth
$h$ in a fluid of density $\rho$. The fluid pushes down on the top with $(P_0+\rho gh)A$ and up on the
bottom with $(P_0+\rho g(h+\ell))A$. The side pushes cancel.

\stepn{2} The net upward force is
\[
F_B=\rho g\ell A=\rho gV,
\]
where $V=\ell A$ is the block's volume. Since $\rho V$ is the mass of fluid that would fill that volume,
$F_B$ is the weight of the displaced fluid. The depth $h$ has cancelled.

\stepn{3} \textbf{Any shape.} Now imagine the body removed and its place taken by fluid of the same shape.
That fluid is at rest, so the pressure forces of the surroundings on its surface must exactly hold up its
weight, $\rho Vg$, and they act along a line through its centre of mass. The surrounding fluid cannot know
what fills the hole, so it pushes on the real body with exactly the same forces. Hence, for a body of any
shape, fully or partly immersed:
\[
\boxed{\;F_B=\text{weight of the fluid displaced}=\rho V_{\text{disp}}\,g\;}
\]
acting through the centre of mass of the displaced fluid (the \textbf{centre of buoyancy}).
\end{derivbox}

\begin{trapbox}{what buoyancy does not depend on}
For a liquid of uniform density, the buoyant force does not depend on how deep the body is, nor on what it
is made of: a lead ball and an aluminium ball of the same volume feel the same buoyant force. Whether the
body sinks depends on how that force compares with its weight, which is where its density enters.
\end{trapbox}

\begin{keybox}{The law of floating}
A body floats at rest when the buoyant force balances its weight. It then sinks until it has displaced its
\emph{own weight} of fluid, so for a body of average density $\rho_{\text{body}}$ in a fluid of density
$\rho_{\text{fluid}}$ the submerged fraction is
\[
\frac{V_{\text{sub}}}{V_{\text{body}}}=\frac{\rho_{\text{body}}}{\rho_{\text{fluid}}}.
\]
A body whose average density is greater than the fluid's cannot displace enough fluid and sinks.
\end{keybox}

\why{Shaped into a hull, a ship's steel encloses so much air that its \emph{average} density is less than
water's, so it floats. Ice ($\rho\approx\SI{917}{kg/m^3}$) floats with about $92\%$ of its volume below the
water, which is why the visible part of an iceberg is so misleading. You float more easily in sea water
than in a pool because sea water is denser. A submarine dives and surfaces by flooding and emptying ballast
tanks, changing its average density; a fish's swim bladder does the same job.}

\begin{exbox}{The apparent weight of a sunken block}\label{ex:fluids-ironblock}
A block of iron of volume \SI{1.0}{\litre} ($\rho_{\text{Fe}}=\SI{7800}{kg/m^3}$) hangs from a spring
balance, fully submerged in water. What does the balance read? ($g=\SI{9.8}{m/s^2}$.)

\stepn{1} True weight: $W=\rho_{\text{Fe}}Vg=7800\times10^{-3}\times9.8\approx\SI{76.4}{N}$.

\stepn{2} Buoyant force $=$ weight of displaced water $=1000\times10^{-3}\times9.8=\SI{9.8}{N}$.

\stepn{3} The balance supports the rest: $W-F_B=76.4-9.8\approx\SI{66.6}{N}$.

\medskip
The block ``loses'' exactly the weight of the water it displaces. The ratio of true weight to loss of
weight, $76.4/9.8=7.8$, is the block's density relative to water, which is how Archimedes could test a crown
without melting it. \ans{$\approx\SI{66.6}{N}$}
\end{exbox}

\begin{exbox}{Does the glass overflow when the ice melts?}\label{ex:fluids-ice}
A glass is filled to the brim with water, with an ice cube floating in it. When the ice melts, does the
water overflow, stay level, or fall?

\stepn{1} While floating, the ice displaces its own weight of water, that is, a volume of water whose mass
equals the mass $m$ of the ice.

\stepn{2} When it melts, the ice becomes water of the same mass $m$, which occupies exactly the volume it was
displacing. The level does not change.

\stepn{3} If the cube had a pebble frozen inside, the floating ice would displace the pebble's weight of
water, a volume larger than the pebble's own. After melting the pebble sinks and displaces only its own
volume, so the level \emph{falls}. Sea ice melting does not raise sea level for the reason in Step 2
(apart from the small difference between salt and fresh water); ice sliding off land into the sea does.
\ans{the level is unchanged}
\end{exbox}

\begin{remarkbox}{The limits of Archimedes' principle}
The principle needs the fluid to press on the body from below. A block pressed flat against the bottom of a
tank, with no water underneath it, feels no upward push at all, only the downward push on its top; this is
how suction cups and stuck ships' hulls hold. The principle also assumes the fluid is at rest in a uniform
field: in a lift accelerating upwards with acceleration $a$ both the body's apparent weight and the buoyant
force scale with $g+a$ (the pseudo-force of \chref{ch:circdyn} acts on the displaced fluid too), so a
floating body floats at the same depth.
\end{remarkbox}

\section{Streamline Flow and the Equation of Continuity}\label{sec:continuity}

Now let the fluid move. A real flow can be fiendishly complicated, so we first build an \textbf{ideal
fluid}, the simplest model that still captures most everyday flows, and say exactly what it assumes.

\begin{defbox}{Steady flow, streamlines and the ideal fluid}
A flow is \textbf{steady} if the velocity of the fluid at each fixed point does not change with time
(different points may have different velocities). In steady flow each particle of fluid follows a fixed
path, a \textbf{streamline}; the velocity at any point is along the tangent to the streamline there, so two
streamlines never cross. A bundle of streamlines forms a \textbf{tube of flow}, whose walls fluid never
crosses. Steady, smooth flow in layers is called \textbf{laminar}; flow with irregular, ever-changing eddies
is \textbf{turbulent}. An \textbf{ideal fluid} is incompressible (constant density) and non-viscous (no
internal friction), and its flow is steady.
\end{defbox}

The first law of an ideal flow is conservation of mass, and it has an immediate, useful consequence.

\begin{derivbox}{The equation of continuity}
\stepn{1} Consider a tube of flow (or a pipe) whose cross-section changes from $A_1$, where the speed is $v_1$,
to $A_2$, where the speed is $v_2$. In a short time $dt$ the fluid at the first section advances
$v_1\,dt$, so a volume $A_1v_1\,dt$ crosses it; similarly $A_2v_2\,dt$ crosses the second.

\stepn{2} No fluid crosses the walls, the flow is steady (nothing accumulates between the sections), and
the density is constant (the fluid cannot be squeezed to fit). So the two volumes are equal:
\[
\boxed{\;A_1v_1=A_2v_2\;}
\]
The product $Av$ is the \textbf{volume flow rate} (\si{m^3/s}), the same at every section.

\stepn{3} Where the pipe narrows, the fluid speeds up; halve the area and the speed doubles.
\end{derivbox}

\inwords{Squeeze the end of a garden hose and the water jets out faster: the same volume each second through
a smaller opening must move faster. A river runs fast through a narrow gorge and slowly across a wide plain
for the same reason.}

\begin{exbox}{Why a stream of water from a tap narrows}\label{ex:fluids-tap}
Water leaves a tap of radius \SI{5.0}{mm} at \SI{1.0}{m/s} and falls freely. Find the radius of the stream
\SI{20}{cm} below the tap ($g=\SI{9.8}{m/s^2}$).

\stepn{1} Each parcel of water is in free fall, so (\chref{ch:straightline}) its speed after falling
$y=\SI{0.20}{m}$ is $v=\sqrt{v_0^2+2gy}=\sqrt{1+3.92}\approx\SI{2.22}{m/s}$.

\stepn{2} Continuity: $\pi r^2v=\pi r_0^2v_0$, so
$r=r_0\sqrt{v_0/v}=5.0\times\sqrt{1.0/2.22}\approx\SI{3.4}{mm}$.

\medskip
The stream thins because it speeds up. Further down it eventually breaks into drops, which is surface
tension at work (Section~\ref{sec:surface-tension}). \ans{$\approx\SI{3.4}{mm}$}
\end{exbox}

\begin{trapbox}{continuity is about volume, which needs constant density}
$A_1v_1=A_2v_2$ conserves \emph{volume}, so it needs an incompressible fluid. Liquids qualify. So, perhaps
surprisingly, does air at speeds well below the speed of sound (below about a third of it, roughly
\SI{100}{m/s}), because the pressure changes in such flows are too small to change its density appreciably.
For a compressible flow the conserved quantity is the mass flow rate, $\rho Av$.
\end{trapbox}

\section{Bernoulli's Equation}\label{sec:bernoulli}

Continuity says a fluid speeds up in a narrow section. Speeding up a parcel of fluid needs a net force on
it, and in a fluid that force can only come from a difference of pressure (and from gravity). So the
pressure must be lower where the fluid moves faster. Making this quantitative is a matter of energy
accounting: apply the work--energy theorem of \chref{ch:energy}, $W_{\text{net}}=\Delta K$, to a slug of
moving fluid (Figure~\ref{fig:fluids-bernoulli}).

\begin{figure}[htbp]
\centering
\tikzfig{fluids_bernoulli}
\figcaption{fig:fluids-bernoulli}{A tube of flow for the work--energy derivation. Follow the fluid between the
two ends. In a short time its back face advances $dx_1$ and its front face $dx_2$, so the net effect is to
move the shaded volume $dV=A_1\,dx_1=A_2\,dx_2$ from the wide, low, slow inlet to the narrow, high, fast
outlet. The fluid behind pushes it forward with $P_1A_1$; the fluid ahead pushes back with $P_2A_2$.}
\end{figure}

\begin{derivbox}{Bernoulli's equation from the work--energy theorem}
\stepn{1} \textbf{Set up the two ends.} The flow is ideal (steady, incompressible, non-viscous). At the inlet
the area is $A_1$, the speed $v_1$, the pressure $P_1$ and the height $h_1$ above a chosen datum; at the
outlet, $A_2$, $v_2$, $P_2$, $h_2$. In a short time the back face of the slug advances $dx_1$ and its front
face $dx_2$, with $dV=A_1\,dx_1=A_2\,dx_2$ by continuity; the mass of that volume is $dm=\rho\,dV$.

\stepn{2} \textbf{Work done by pressure.} The fluid behind pushes the slug forward with force $P_1A_1$ through
$dx_1$, doing work $P_1A_1\,dx_1=P_1\,dV$. The fluid ahead pushes back with $P_2A_2$ through $dx_2$, doing
negative work $-P_2\,dV$. The walls push perpendicular to the flow and do no work. Net pressure work:
\[
W_{\text{pressure}}=(P_1-P_2)\,dV.
\]

\stepn{3} \textbf{Work done by gravity.} Between the two instants, the fluid in the middle of the tube is
unchanged (the flow is steady); the net effect is that the mass $dm$ has moved from height $h_1$ to $h_2$.
Gravity therefore does work
\[
W_{\text{gravity}}=-dm\,g\,(h_2-h_1)=-\rho\,dV\,g\,(h_2-h_1).
\]

\stepn{4} \textbf{Change of kinetic energy.} By the same argument, the kinetic energy changes as if $dm$ had
gone from speed $v_1$ to $v_2$: $\Delta K=\tfrac12\rho\,dV\,(v_2^2-v_1^2)$. With no viscous friction there is
no other work, and $W_{\text{pressure}}+W_{\text{gravity}}=\Delta K$ reads
\[
(P_1-P_2)\,dV-\rho\,dV\,g\,(h_2-h_1)=\tfrac12\rho\,dV\,(v_2^2-v_1^2).
\]

\stepn{5} \textbf{Divide by $dV$ and collect the ends.}
\[
\boxed{\;P_1+\tfrac12\rho v_1^2+\rho gh_1=P_2+\tfrac12\rho v_2^2+\rho gh_2\;}
\]
The two ends were arbitrary, so the combination has the same value all along the tube of flow.
\end{derivbox}

\begin{keybox}{Bernoulli's equation}
Along a streamline of a steady flow of an incompressible, non-viscous fluid,
\[
P+\tfrac12\rho v^2+\rho gh=\text{constant}.
\]
Each term is an energy per unit volume (\si{J/m^3}, the same unit as \si{Pa}): $\tfrac12\rho v^2$ is the
kinetic energy per unit volume, $\rho gh$ the gravitational potential energy per unit volume, and $P$ the
work per unit volume that the pressure can do in pushing the fluid along. For fluid at rest ($v=0$) the
equation reduces to $P+\rho gh=\text{constant}$, the depth law of Section~\ref{sec:fluid-pressure}.
\end{keybox}

\inwords{Bernoulli's equation is energy conservation for a flowing fluid. If a fluid speeds up at the same
height, its pressure must drop to pay for the extra kinetic energy; if it rises at the same speed, its
pressure drops to pay for the extra potential energy. Speed and height are bought with pressure.}

\begin{trapbox}{low pressure is not suction}
``Fast flow means low pressure'' does not mean the moving fluid pulls on anything. Pressure is always a
push. A region of fast flow simply pushes \emph{less} than the slower fluid around it, and the
\emph{difference} of the pushes is the net force. The cause and effect also run the other way: in a pipe
it is the pressure difference that accelerates the fluid into the narrow section.
\end{trapbox}

\subsection{Torricelli's law}

The first application asks how fast water leaves a hole in a tank (Figure~\ref{fig:fluids-torricelli}).

\begin{figure}[htbp]
\centering
\tikzfig{fluids_torricelli}
\figcaption{fig:fluids-torricelli}{Torricelli's law. Point 1 on the free surface and point 2 in the jet just
outside the hole lie on one streamline (dashed). Both are open to the air, so both are at pressure $P_0$;
the surface barely moves, so $v_1\approx0$. Bernoulli's equation turns the height $h$ into the jet speed
$v=\sqrt{2gh}$, after which the water is a projectile and lands a distance $2\sqrt{hH}$ from the tank.}
\end{figure}

\begin{derivbox}{Torricelli's law: the speed of efflux}
A wide tank of water has a small hole a depth $h$ below the free surface. Apply Bernoulli's equation between
two well-chosen points on one streamline and evaluate its three terms at each.

\stepn{1} \textbf{Choose the points.} Point~1 is on the free surface; point~2 is in the jet just outside the
hole. Then
\[
P_1+\tfrac12\rho v_1^2+\rho gh_1=P_2+\tfrac12\rho v_2^2+\rho gh_2.
\]

\stepn{2} \textbf{Pressure terms.} The free surface and the emerging jet are both open to the atmosphere,
so $P_1=P_2=P_0$, and these terms cancel. (The jet is a free stream surrounded by air; nothing confines it
to push it above atmospheric pressure.)

\stepn{3} \textbf{Kinetic terms.} By continuity, the surface falls at a speed smaller than the jet's by the
ratio of the hole's area to the tank's, so for a wide tank $v_1\approx0$. At the hole the water leaves
with the speed $v$ we want.

\stepn{4} \textbf{Potential terms.} Measuring heights from the hole, $h_1=h$ and $h_2=0$.

\stepn{5} \textbf{Combine.}
\[
\cancel{P_0}+0+\rho gh=\cancel{P_0}+\tfrac12\rho v^2+0
\quad\Longrightarrow\quad\boxed{\;v=\sqrt{2gh}\;}
\]
The density has cancelled too: a denser liquid has more potential and more kinetic energy per unit volume
in the same proportion.

\medskip
This is exactly the speed of a body falling freely from rest through $h$ (\chref{ch:straightline}),
because the physics is the same: potential energy becoming kinetic energy. After leaving the hole
horizontally, the jet is a projectile (\chref{ch:projectiles}). From a hole at height $H$ above the floor
it falls for $t=\sqrt{2H/g}$ and lands a horizontal distance $vt=2\sqrt{hH}$ away. For a tank filled to a
fixed depth $D=h+H$, this range is largest when $h=H=D/2$ (the product $hH$ of two numbers with a fixed
sum is largest when they are equal), and the jet then lands a distance $D$ away.

\medskip
\noindent\textbf{Two ways to see the $\rho gh$.} A natural worry: should the pressure at the hole not be
$P_0+\rho gh$, since the hole is at depth $h$? Inside the tank, just beside the hole, it is. There are two
valid ways to keep the books, and they agree.
\begin{itemize}
\item \textbf{Surface to jet (used above).} The $\rho gh$ appears as the \emph{potential energy} of the high
free surface. Both ends are at $P_0$.
\item \textbf{Still water at the hole's depth to jet.} Start instead from a point inside the tank, at the
hole's level but far enough from it that the water is nearly still. There the pressure really is
$P_0+\rho gh$ and the height is zero, and Bernoulli gives $P_0+\rho gh=P_0+\tfrac12\rho v^2$: the same
$v=\sqrt{2gh}$. Now $\rho gh$ appears as pressure.
\end{itemize}
Both work because hydrostatic pressure \emph{is} the weight of the column above: the same $\rho gh$ can be
charged to height at the top or to pressure at the bottom. What you must not do is count it twice, by
taking the surface point with its height and also putting $P_0+\rho gh$ in the jet; that gives a spurious
$v=\sqrt{4gh}$.
\end{derivbox}

\begin{remarkbox}{How good is $\sqrt{2gh}$?}
Real jets come out a few per cent slower, because a little energy is lost to viscosity at the edge of the
hole. A sharp-edged hole also produces a jet that contracts just outside it to about $60\%$ of the hole's
area (the \emph{vena contracta}), because the streamlines converging on the hole cannot turn sharply.
So $\sqrt{2gh}$ gives the speed well, but the volume flow rate is usually about $0.6$ of $A\sqrt{2gh}$.
\end{remarkbox}

\subsection{The Venturi meter}

The second application measures a flow from the pressure drop it causes in a constriction, which is the
continuity equation and Bernoulli's equation working together.

\begin{exbox}{A Venturi meter}\label{ex:fluids-venturi}
Water flows through a horizontal pipe of diameter \SI{10}{cm}, which narrows smoothly to a throat of
diameter \SI{5.0}{cm}. Pressure tappings show that the pressure in the throat is \SI{15}{kPa} lower than in
the wide pipe. Find the speed in the wide pipe and the volume flow rate.

\stepn{1} \textbf{Continuity.} Halving the diameter quarters the area, so $v_2=4v_1$.

\stepn{2} \textbf{Bernoulli at one height.} $P_1-P_2=\tfrac12\rho(v_2^2-v_1^2)=\tfrac12\rho(16-1)v_1^2$, so
\[
v_1=\sqrt{\frac{2(P_1-P_2)}{15\rho}}=\sqrt{\frac{2(15000)}{15(1000)}}=\sqrt2\approx\SI{1.41}{m/s}.
\]

\stepn{3} \textbf{Flow rate.} $A_1v_1=\pi(0.050)^2(1.41)\approx\SI{1.11e-2}{m^3/s}$, about \SI{11}{\litre}
each second.

\medskip
The same physics, a fast stream with lowered pressure, draws liquid up a tube into the air stream of a
perfume atomiser or an old car's carburettor. \ans{\SI{1.41}{m/s}; \SI{11}{L/s}}
\end{exbox}

\why{Bernoulli's equation explains a wide range of effects. Two ships moving side by side are pushed
together, because the water squeezed between them flows faster and presses less. A strong wind over a
roof lowers the pressure above it while the still air inside presses as before, and roofs are lifted off.
Wind blowing across a chimney top lowers the pressure there and helps draw the smoke up. A spinning football
drags air round with it, so the air flows faster on one side than the other, and the pressure difference
bends its path (the Magnus effect).}

\begin{remarkbox}{What Bernoulli does and does not explain about lift}
Over an aircraft wing the air flows faster above than below, and the pressure above is correspondingly
lower; Bernoulli's equation relates the two correctly, and the pressure difference times the wing area is
the lift. What Bernoulli does not explain is \emph{why} the air goes faster over the top. A popular
explanation says that air parcels split at the front of the wing must ``meet again'' at the back, so the
air over the longer upper surface must go faster. This is false: wind-tunnel pictures show the upper air
arriving at the back well \emph{before} the lower air, and nothing requires the parcels to meet. The full
answer is that the wing, by its shape and its tilt (angle of attack), turns the oncoming air
\emph{downwards}. By Newton's third law (\chref{ch:newton}) the air pushes the wing up, and the pressure
pattern round the wing is how that push is delivered. This is why a flat paper plane flies, and why an
aerobatic aircraft can fly upside down by tilting its wing. Lift as ``pressure difference'' and lift as
``reaction to the downward-turned air'' are the same force counted two ways.
\end{remarkbox}

\begin{remarkbox}{Estimate before you believe: the bottle ``eco-cooler''}
A widely shared idea from rural Bangladesh mounts cut plastic bottles through a board in a window, wide end
outside and narrow neck inside, and claims that the hut is cooled by several degrees with no electricity.
Two explanations circulate. Both can be checked with this chapter, and both fail.

\smallskip
\textbf{``The bottle compresses the air, which expands and cools.''} The largest pressure a breeze of speed
$v=\SI{2}{m/s}$ can build up against a board is about its kinetic energy per unit volume,
$\tfrac12\rho v^2=\tfrac12(1.2)(2)^2\approx\SI{2.4}{Pa}$ (Bernoulli, bringing the air to rest). Even if all
of that energy per unit volume were taken out of the air as heat, the temperature would fall by
$\Delta P/(\rho s_p)$, where $s_p\approx\SI{1005}{J.kg^{-1}.K^{-1}}$ is air's specific heat capacity at constant pressure (heat per kilogram per kelvin, defined in \chref{ch:thermal} and predicted from molecular structure in \chref{ch:kinetic}): $2.4/(1.2\times1005)\approx\SI{0.002}{K}$. That is two thousandths of a degree.

\smallskip
\textbf{``The funnel speeds the breeze up, by continuity.''} Continuity applies only to air that actually
passes through the bottle, and the speed of that air is set by the pressure difference across the board.
By Bernoulli, \SI{2.4}{Pa} can drive air through any opening at no more than $\sqrt{2\Delta P/\rho}$, which
is about the wind's own speed of \SI{2}{m/s}. The funnel cannot make the jet faster than the wind; most
of the air arriving at the wide mouth simply spills round the board, and the bottles let \emph{less} air
in than an open window would.

\smallskip
So the device cannot cool the incoming air, and careful measurements have not found the large
temperature drops claimed for it. It can do no more than any opening does: let in a draught, which feels
cooler on the skin because moving air carries heat and sweat away faster. The lesson is general: estimate
the size of an effect before believing the story told about it.
\end{remarkbox}

\begin{remarkbox}{Where Bernoulli's equation stops working}
The derivation assumed a steady flow of an incompressible, non-viscous fluid, and it applies along one
streamline. Viscosity is the commonest failure: water pumped through a long horizontal pipe of uniform
width keeps the same speed (continuity), yet its pressure falls steadily along the pipe, because work is
being done against internal friction and turned into heat. Turbulence breaks the steady-flow assumption,
and at speeds approaching the speed of sound air can no longer be treated as incompressible. Within these
limits the equation is extremely accurate; beyond them it still gives the right trends, but not the
numbers.
\end{remarkbox}

\section{Viscosity}\label{sec:viscosity}

The ideal fluid has no internal friction, so nothing in it would ever slow down. Real fluids do: stirred
tea stops swirling, and honey poured from a spoon drags behind it. The internal friction of a fluid is its
\textbf{viscosity}. Its effects are usually small for water and air in pipes of everyday size, but they
dominate the motion of small bodies, thin films and slow flows, and they are why Bernoulli's equation fails
along long pipes.

\begin{figure}[htbp]
\centering
\tikzfig{fluids_viscosity}
\figcaption{fig:fluids-viscosity}{Viscous flow between a fixed plate and a moving one. The liquid sticks to
both plates, so its speed rises from zero at the bottom to $v_0$ at the top, and each layer slides over the
one below. Keeping the top plate moving at constant speed needs a steady force $F$, which is proportional
to the area and to the velocity gradient $dv/dy$.}
\end{figure}

\begin{defbox}{Coefficient of viscosity}
When layers of a fluid slide past one another (Figure~\ref{fig:fluids-viscosity}), each layer drags its
neighbours. If the speed changes by $dv$ across a thickness $dy$ perpendicular to the flow, the tangential
force needed on an area $A$ of a layer is
\[
F=\eta A\,\frac{dv}{dy}\qquad\text{(Newton's law of viscosity)},
\]
which defines the \textbf{coefficient of viscosity} $\eta$. Its SI unit is \si{N.s/m^2}$=\si{Pa.s}$, and its
dimensions are $\mathsf{M\,L^{-1}\,T^{-1}}$, as found in \chref{ch:measurement}. Typical values near room
temperature: air \SI{1.8e-5}{Pa.s}, water \SI{1.0e-3}{Pa.s}, engine oil about \SI{0.1}{Pa.s}, glycerine
\SI{1.5}{Pa.s}, honey a few \si{Pa.s}.
\end{defbox}

The ratio $F/A$ is a shear stress, the quantity \chref{ch:solids} used for the shearing of a solid. The
contrast is the point: in a solid, shear stress is proportional to the \emph{strain} (how far it has
deformed), and the solid stops deforming; in a fluid, it is proportional to the \emph{rate} of strain
($dv/dy$), and the fluid keeps flowing for as long as the stress acts.

\begin{trapbox}{viscosity changes with temperature, in opposite ways}
The viscosity of a liquid falls steeply as it is heated (warm honey pours easily; water's viscosity halves
between \SI{20}{\celsius} and \SI{60}{\celsius}), because its molecules escape more easily from their
neighbours' attraction. A gas becomes \emph{more} viscous when heated, because its viscosity comes from
molecules carrying momentum between layers, and faster molecules carry more; the mean free path of \chref{ch:kinetic} explains why
(Section~\ref{sec:ktmfp}).
\end{trapbox}

Whether a flow is smooth or turbulent is decided by a competition between the fluid's inertia, which tends
to throw it about, and its viscosity, which smooths it out. For flow at speed $v$ through a pipe of
diameter $D$, the ratio is measured by the dimensionless \textbf{Reynolds number}
\[
Re=\frac{\rho vD}{\eta}.
\]
Experiment shows that flow in a pipe is laminar when $Re$ is below about $2000$ and turbulent above about
$3000$. Water at \SI{0.2}{m/s} in a \SI{1}{cm} pipe has $Re=(1000)(0.2)(0.01)/(10^{-3})=2000$, right at the
boundary, which is why a tap opened gently gives a clear, smooth stream that turns ragged as it is opened
further.

\subsection{Stokes' law and terminal velocity}

A body moving through a fluid drags the neighbouring layers with it, and their viscosity resists: this is
\textbf{viscous drag}. For a small sphere moving slowly, so that the flow round it is smooth, the form of the
drag follows from dimensions alone.

\begin{derivbox}{The form of Stokes' law from dimensions}
\stepn{1} \textbf{What can the drag depend on?} When the flow is smooth and slow, the fluid slides round the
sphere rather than being flung aside, so its inertia (and hence its density) does not matter. The drag $F$
can then depend only on the viscosity $\eta$, the sphere's radius $r$ and its speed $v$. Try
$F=\beta\,\eta^a r^b v^c$ with $\beta$ a pure number.

\stepn{2} \textbf{Match dimensions} (\chref{ch:measurement}):
$\mathsf{M\,L\,T^{-2}}=(\mathsf{M\,L^{-1}\,T^{-1}})^a\,\mathsf{L}^b\,(\mathsf{L\,T^{-1}})^c$. Mass gives
$a=1$; time gives $a+c=2$, so $c=1$; length gives $-a+b+c=1$, so $b=1$. Hence
\[
F=\beta\,\eta rv.
\]
Each factor makes sense: a more viscous fluid, a bigger sphere and a faster motion all shear the fluid
harder.

\stepn{3} \textbf{The number.} Dimensions cannot give $\beta$. A full solution of the flow round the sphere
(beyond this book, so stated without proof) gives $\beta=6\pi$:
\[
\boxed{\;F=6\pi\eta rv\;}\qquad\text{(Stokes' law).}
\]
The feature to carry forward is that this drag is proportional to the speed.
\end{derivbox}

That proportionality produces a terminal velocity. Consider a small sphere of radius $r$ and density
$\rho_s$ released from rest in a fluid of density $\rho<\rho_s$ and viscosity $\eta$.

\begin{derivbox}{Terminal velocity of a falling sphere}
\stepn{1} \textbf{The forces.} Weight $W=\tfrac43\pi r^3\rho_sg$ (down); buoyancy $F_B=\tfrac43\pi r^3\rho g$
(up, Archimedes); drag $6\pi\eta rv$ (up, opposing the motion, and growing with $v$). With
$m=\tfrac43\pi r^3\rho_s$, Newton's second law gives
\[
m\frac{dv}{dt}=\tfrac43\pi r^3(\rho_s-\rho)g-6\pi\eta rv.
\]

\stepn{2} \textbf{At release.} With $v=0$ there is no drag, and the acceleration is greatest:
\[
a_0=\frac{\tfrac43\pi r^3(\rho_s-\rho)g}{\tfrac43\pi r^3\rho_s}=\Big(1-\frac{\rho}{\rho_s}\Big)g.
\]

\stepn{3} \textbf{Terminal velocity.} As $v$ grows, the drag grows and the acceleration shrinks. It reaches
zero when the drag balances weight minus buoyancy; the speed then stops changing. Setting $dv/dt=0$:
\[
6\pi\eta rv_t=\tfrac43\pi r^3(\rho_s-\rho)g\quad\Longrightarrow\quad
\boxed{\;v_t=\frac{2r^2(\rho_s-\rho)g}{9\eta}\;}
\]

\stepn{4} \textbf{How the speed approaches $v_t$.} The equation of Step 1 can be written
$dv/dt=(v_t-v)/\tau$, with $\tau=m/(6\pi\eta r)=2r^2\rho_s/(9\eta)$. The gap $v_t-v$ therefore decays at
a rate proportional to itself, a first-order decay (\chref{ch:toolkit}, Section~\ref{sec:explog}), and
\[
v(t)=v_t\big(1-e^{-t/\tau}\big).
\]
The initial slope is $v_t/\tau=a_0$, as Step 2 requires. After one time constant $\tau$ the sphere has
reached $63\%$ of $v_t$, and after $5\tau$ more than $99\%$ (Figure~\ref{fig:fluids-terminal}).

\stepn{5} \textbf{Why the terminal state is stable.} The sphere has not ``run out of push'': the forces still
act, but they cancel, so by Newton's first law it moves at constant $v_t$. If it were nudged faster, the
drag would exceed the driving force and slow it back; if slower, the driving force would win.
\end{derivbox}

\begin{figure}[htbp]
\centering
\tikzfig{fluids_terminal}
\figcaption{fig:fluids-terminal}{The approach to terminal velocity. The speed rises steeply at first, with
slope $a_0=(1-\rho/\rho_s)g$, because the drag starts at zero; the growing drag then eats into the
acceleration, and the speed levels off at $v_t$, where the forces balance. The tangent at the start would
reach $v_t$ after one time constant $\tau$.}
\end{figure}

The $r^2$ in $v_t$ is the striking part: a droplet twice as wide falls four times as fast. This is why
mist and fog hang in the air while rain falls, and it is the principle behind Millikan's measurement of the
electron's charge from the fall of tiny oil drops.

\begin{exbox}{A fog droplet}\label{ex:fluids-fog}
A water droplet of radius \SI{10}{\micro m} falls through air ($\eta=\SI{1.8e-5}{Pa.s}$,
$\rho_{\text{air}}=\SI{1.2}{kg/m^3}$, negligible beside water's \SI{1000}{kg/m^3}). Find its terminal
velocity and the time it takes to reach it, and check that Stokes' law applies.

\stepn{1} $v_t=\dfrac{2(10^{-5})^2(1000)(9.8)}{9(1.8\times10^{-5})}\approx\SI{1.2e-2}{m/s}$, about a
centimetre per second.

\stepn{2} $\tau=\dfrac{2r^2\rho_s}{9\eta}=\dfrac{2(10^{-5})^2(1000)}{9(1.8\times10^{-5})}\approx\SI{1.2e-3}{s}$.
The droplet reaches its terminal velocity within a few milliseconds, and drifts down at it ever after.

\stepn{3} \textbf{Check.} Using the droplet's diameter as the length,
$Re=\rho_{\text{air}}v_t(2r)/\eta=(1.2)(0.012)(2\times10^{-5})/(1.8\times10^{-5})\approx0.016$, far below~$1$,
so the flow round it is smooth and Stokes' law holds. \ans{\SI{1.2}{cm/s}; about \SI{1}{ms}}
\end{exbox}

\begin{trapbox}{Stokes' law is for small, slow bodies}
Stokes' law assumes smooth flow round the sphere, which requires a Reynolds number (based on the sphere's
diameter) well below~$1$. A raindrop, a falling ball or a skydiver is far outside that range. Their drag
grows roughly as $v^2$, not $v$, and the Stokes formula overestimates their terminal speeds absurdly (see
the last challenge problem of this chapter). Always check $Re$ before using Stokes' law.
\end{trapbox}

\section{Surface Tension}\label{sec:surface-tension}

The second property the ideal fluid leaves out lives at its surface. A steel needle laid gently on water
floats, insects walk on ponds, a drop of water on a waxed surface pulls itself into a bead, and water climbs
up a thin glass tube by itself. All of these say that a liquid's surface behaves like a stretched elastic
skin.

The reason is molecular. A molecule deep inside a liquid is attracted by neighbours on all sides, and the
pulls cancel. A molecule at the surface has neighbours only below and beside it, so it is pulled inwards,
and bringing a molecule from the bulk to the surface costs energy. A liquid therefore has extra energy in
proportion to its surface area, and, like any system, it moves towards lower energy by making its surface
as small as it can.

\begin{figure}[htbp]
\centering
\tikzfig{fluids_wireframe}
\figcaption{fig:fluids-wireframe}{Measuring surface tension with a soap film. The film on the frame has two
surfaces, front and back, each pulling the sliding wire inwards; holding the wire still needs $F=2S\ell$.
Pulling it out by $dx$ does work $F\,dx$ and creates new surface of area $2\ell\,dx$, so the work per unit
of new area is $S$.}
\end{figure}

\begin{defbox}{Surface tension and surface energy}
Draw a line of length $\ell$ on a liquid surface. The surface on each side pulls the line towards itself,
along the surface and perpendicular to the line, with a force proportional to $\ell$. The \textbf{surface
tension} $S$ is that force per unit length:
\[
S=\frac{F}{\ell}\qquad(\si{N/m}).
\]
It acts along the surface, never out of it, and always so as to shrink the surface. Enlarging the surface
by an area $\Delta A$ at constant temperature takes work $W=S\,\Delta A$
(Figure~\ref{fig:fluids-wireframe}), so $S$ is equally the \textbf{surface energy} per unit area
(\si{J/m^2}$=$\si{N/m}). Typical values near \SI{20}{\celsius}: water \SI{0.072}{N/m}, soap solution
about \SI{0.03}{N/m}, mercury about \SI{0.47}{N/m}.
\end{defbox}

\inwords{Surface tension and surface energy are one idea seen two ways: a force pulling along the surface,
or an energy price for having surface at all. A force of $S$ per metre, pulled through a distance, creates
$S$ joules of surface energy per square metre.}

This book writes surface tension as $S$, following NCERT; many books write $T$ or $\gamma$, but in this
volume $T$ is always temperature and $\gamma$ the ratio of heat capacities. Surface tension falls as the
temperature rises, and detergents lower it sharply, which is how soap lets water spread over and wet a
greasy surface it would otherwise bead up on.

\begin{trapbox}{a soap film has two surfaces}
A soap film is a thin sheet of liquid with air on both sides, so it has two surfaces. Every force or energy
for a film carries a factor $2$: the frame needs $F=2S\ell$, and blowing a soap bubble of radius $R$
creates surface area $2\times4\pi R^2$. A liquid drop has only one surface.
\end{trapbox}

\subsection{The angle of contact}

Where a liquid surface meets a solid wall, the shape of the surface depends on a competition: the liquid's
molecules attract each other (cohesion) and are attracted by the wall (adhesion). The outcome is recorded
by one angle (Figure~\ref{fig:fluids-contactangle}).

\begin{figure}[htbp]
\centering
\tikzfig{fluids_contactangle}
\figcaption{fig:fluids-contactangle}{The angle of contact $\theta$ is measured from the wall, inside the
liquid, to the tangent to the liquid surface at the contact point. A wetting liquid (left) climbs the wall,
its surface curving up to meet it, and $\theta$ is small; a non-wetting liquid (right) is pushed down at the
wall, and $\theta$ exceeds $90^\circ$. Measuring to the normal instead of the tangent would swap every
$\cos\theta$ for a $\sin\theta$.}
\end{figure}

\begin{defbox}{Angle of contact}
At the line where a liquid surface meets a solid, draw the tangent to the liquid surface. The \textbf{angle
of contact} $\theta$ is the angle between that tangent and the solid surface, measured through the liquid.
If $\theta<90^\circ$ the liquid \textbf{wets} the solid (adhesion beats cohesion) and its surface curves up
against the wall: water on clean glass, $\theta\approx0$. If $\theta>90^\circ$ it does not wet the solid
and its surface curves down: mercury on glass, $\theta\approx140^\circ$; water on wax, about $105^\circ$.
\end{defbox}

The tangent, not the normal, is used because surface tension pulls along the surface: the tangent's
direction decides how much of that pull points along the wall.

\subsection{Why a free drop is a sphere}

A liquid makes its surface energy as small as it can. For a drop with no other forces acting, that fixes
its shape.

\begin{derivbox}{Least area at fixed volume}
\stepn{1} \textbf{The question.} A free drop has a fixed volume $V$ (the liquid is nearly incompressible)
and surface energy $S\times(\text{area})$. It settles into the shape of least energy, hence of least area
for its volume. So the question is purely geometric: which closed shape of given volume has the least
surface area?

\stepn{2} \textbf{The sphere.} For a sphere of radius $R$, $V=\tfrac43\pi R^3$ and $A=4\pi R^2$. Eliminating
$R$,
\[
A_{\text{sphere}}=4\pi\left(\frac{3V}{4\pi}\right)^{2/3}=(36\pi)^{1/3}V^{2/3}\approx4.84\,V^{2/3}.
\]

\stepn{3} \textbf{A cube for comparison.} A cube of the same volume has side $V^{1/3}$ and area
$6V^{2/3}$, about $24\%$ more than the sphere. Every other shape tried (cylinder, ellipsoid, teardrop) also
does worse.

\stepn{4} \textbf{The general statement.} The \emph{isoperimetric inequality} of geometry (stated here
without proof, because the proof needs the calculus of variations) says that every closed surface satisfies
$A^3\ge36\pi V^2$, with equality only for the sphere. So no shape of volume $V$ has less area than the
sphere, and since $S>0$, the sphere is the unique shape of least surface energy.

\stepn{5} \textbf{Why only free drops are round.} A drop resting on a table or hanging from a tap also has
gravitational potential energy, and its real shape minimises the sum. Surface energy scales as $R^2$ and
gravitational energy as $R^4$ (mass $\propto R^3$ times height $\propto R$), so small drops are nearly
spherical and large ones flatten into puddles. Water floating in an orbiting spacecraft is round at any
size.
\end{derivbox}

\subsection{Excess pressure inside a drop and a bubble}

A curved surface under tension squeezes what is inside it, just as a stretched balloon skin does. So the
pressure just inside a curved liquid surface exceeds the pressure just outside, by an amount we call the
\textbf{excess pressure} $\Delta P=P_{\text{in}}-P_{\text{out}}$ (Figure~\ref{fig:fluids-drop}).

\begin{figure}[htbp]
\centering
\tikzfig{fluids_drop}
\figcaption{fig:fluids-drop}{Excess pressure in a drop, by a force balance. Cut the drop through its centre and
look at the right half. The excess pressure pushes it to the right with $\Delta P$ times the area of the
cut, $\pi R^2$ (crimson); the surface of the left half pulls its rim, of length $2\pi R$, to the left with
$S$ per unit length (orange, an applied pull told apart from the pressure force). Equating the two gives
$\Delta P=2S/R$.}
\end{figure}

\begin{derivbox}{Excess pressure inside a drop and a bubble}
\stepn{1} \textbf{Force balance on half a drop.} In Figure~\ref{fig:fluids-drop} the right half of a drop of
radius $R$ is pushed to the right by the liquid of the left half, with pressure $P_{\text{in}}$ on the
circular cut of area $\pi R^2$. The outside air pushes on its curved surface; the net effect of a uniform
pressure on a curved surface is the pressure times the area projected onto the cut, so the air's net push
is $P_{\text{out}}\pi R^2$ to the left. The surface of the left half pulls along the rim with
$S\times2\pi R$, to the left. Balance:
\[
\Delta P\,\pi R^2=2\pi RS\quad\Longrightarrow\quad\boxed{\;\Delta P_{\text{drop}}=\frac{2S}{R}\;}
\]

\stepn{2} \textbf{The same result from energy.} Let the drop grow by $dR$. The excess pressure does work
$\Delta P\times(4\pi R^2)\,dR$ pushing the surface out, and the surface energy rises by
$S\,d(4\pi R^2)=8\pi SR\,dR$. In equilibrium the two are equal, which again gives $\Delta P=2S/R$.

\stepn{3} \textbf{A soap bubble has two surfaces}, inner and outer, each of nearly the same radius. Each
contributes $2S/R$, so
\[
\boxed{\;\Delta P_{\text{bubble}}=\frac{4S}{R}\;}
\]
\stepn{4} \textbf{Read it.} The smaller the radius, the larger the excess pressure. That is why the first
puff is the hardest when blowing up a bubble (or a balloon), and why, when a small and a large soap bubble
are connected by a tube, the small one empties into the large one.
\end{derivbox}

\subsection{Capillary rise}

The excess-pressure law and the angle of contact together explain why a liquid climbs or sinks in a narrow
tube, a \textbf{capillary} (Latin \emph{capillus}, a hair), as in Figure~\ref{fig:fluids-capillary}.

\begin{figure}[htbp]
\centering
\tikzfig{fluids_capillary}
\figcaption{fig:fluids-capillary}{Capillarity, drawn to scale for equal tubes. Water, which wets glass, climbs
the tube and forms a concave meniscus; mercury, which does not, is pushed down below the outside level
and forms a convex one. The mercury's depression is only about a third of the water's rise, because its
much larger density outweighs its larger surface tension.}
\end{figure}

\begin{figure}[htbp]
\centering
\tikzfig{fluids_capfbd}
\figcaption{fig:fluids-capfbd}{The forces that hold up a capillary column. At the contact line the liquid
surface pulls on the glass along its own tangent (dashed); by Newton's third law the glass pulls the liquid
back with $S$ per unit length the other way. Its vertical part, $S\cos\theta$, lifts the column; its
horizontal part, $S\sin\theta$, cancels round the circular rim.}
\end{figure}

\begin{derivbox}{Capillary rise: $h=2S\cos\theta/(\rho gr)$}
\stepn{1} \textbf{Where the pull acts.} The liquid meets the glass along a circle at the top of the column, the
contact line, of length $2\pi r$, where $r$ is the tube's inner radius. Along it the liquid surface pulls on
the glass with $S$ per unit length, directed along the surface's tangent, which leaves the wall at the
angle of contact $\theta$ (Figure~\ref{fig:fluids-capfbd}).

\stepn{2} \textbf{The pull on the liquid.} By Newton's third law the glass pulls the liquid with an equal and
opposite force, $S$ per unit length. Its component along the wall, $S\cos\theta$, points up; its component
across the wall, $S\sin\theta$, points towards the glass, and round the circular rim these horizontal pulls
cancel. The total upward force on the liquid is
\[
F_{\text{up}}=2\pi rS\cos\theta.
\]
A strongly wetting liquid ($\theta\to0$) pulls almost straight up; at $\theta=90^\circ$ there is no vertical
pull and no rise.

\stepn{3} \textbf{The weight to be held.} The raised liquid is very nearly a cylinder of radius $r$ and height
$h$ (measured to the bottom of the meniscus), with weight $\rho\pi r^2hg$. The small extra liquid around
the meniscus adds a volume $\pi r^3/3$ for a hemispherical meniscus, which replaces $h$ by $h+r/3$; for a
tube of radius \SI{0.2}{mm}, where $h\approx\SI{7}{cm}$, that is a correction of about $0.1\%$, and it is
neglected.

\stepn{4} \textbf{Balance.} The pressure at the foot of the column equals the pressure at the same level in
the open liquid outside, so the outside liquid neither helps nor hinders, and
\[
2\pi rS\cos\theta=\rho\pi r^2hg\quad\Longrightarrow\quad\boxed{\;h=\frac{2S\cos\theta}{\rho gr}\;}
\]

\stepn{5} \textbf{Check by pressure.} The meniscus is a spherical cap meeting the wall at angle $\theta$, so
its radius of curvature is $R=r/\cos\theta$, and it is concave upwards for a wetting liquid. The pressure
just below it is therefore lower than the air pressure above by $2S/R=2S\cos\theta/r$. But by the depth
law it is also $P_0-\rho gh$, the pressure a height $h$ above the outside level. Equating the two gives the
same $h$.

\stepn{6} \textbf{Read it.} The rise is inversely proportional to the radius: halve the tube and the liquid
climbs twice as high. For $\theta>90^\circ$, $\cos\theta<0$ and $h<0$: mercury is \emph{depressed}.
\end{derivbox}

\begin{trapbox}{a capillary never overflows}
If a tube is shorter than the height $h$ the liquid would reach, the liquid does not spill out of the top.
It rises to the rim and stops there, with a flatter meniscus (larger radius of curvature) whose smaller
excess pressure exactly supports the shorter column. Capillarity cannot make a perpetual fountain.
\end{trapbox}

\begin{remarkbox}{Where the capillary formula stops working}
The derivation treats the meniscus as a spherical cap, which is right only when gravity cannot distort it,
that is, when the tube's radius is small compared with the \emph{capillary length}
$\sqrt{S/(\rho g)}$, about \SI{2.7}{mm} for water. In a wide tube the surface is flat in the middle and the
rise is far less than $2S\cos\theta/(\rho gr)$ predicts. The formula also needs a clean tube: a trace of
grease raises $\theta$ and cuts the rise sharply.
\end{remarkbox}

\why{Viscosity is why honey pours slowly, why engine oil is graded by thickness, and why a narrowed blood
vessel needs a higher pressure to keep the same flow. Surface tension is why raindrops are round, why soap
cleans, why a paper towel soaks up a spill, and why a needle can float. Both are small forces that are easy
to ignore, and both are indispensable to living things and to technology.}

\begin{remarkbox}{How a redwood drinks}
How does a \SI{100}{m} tree lift water from its roots to its highest leaves? The tempting answer, capillary
rise in the trunk, fails, and seeing why uses several ideas of this section.

\smallskip
\textbf{Capillarity alone is far too weak.} A xylem vessel (the tree's water pipe) has a radius of about
$r=\SI{20}{\micro m}$. For water wetting the wall,
\[
h=\frac{2S}{\rho gr}=\frac{2(0.072)}{(1000)(9.8)(20\times10^{-6})}\approx\SI{0.7}{m}.
\]
Enough for a seedling; useless for a redwood.

\smallskip
\textbf{The engine is evaporation at the leaves.} Water fills the microscopic pores of the cell walls inside
a leaf, and as it evaporates, the surfaces retreat into pores only a few nanometres across. By the
excess-pressure law, a meniscus of radius $R\approx\SI{5}{nm}$ can sustain a pressure difference of up to
\[
\Delta P=\frac{2S}{R}=\frac{2(0.072)}{5\times10^{-9}}\approx\SI{2.9e7}{Pa}\approx\SI{280}{atm}.
\]
Holding up a \SI{100}{m} column needs $\rho gh\approx\SI{9.8e5}{Pa}$, about \SI{10}{atm}, plus roughly as
much again to overcome the viscous resistance of the flow. The menisci can pull far harder than this; in
practice the suction measured in tall trees is a few tens of atmospheres at most.

\smallskip
\textbf{Cohesion carries the pull down.} Water molecules attract one another strongly (the same cohesion
that gives water its high surface tension), so the water in the xylem hangs together as an unbroken thread
from leaf to root, under \emph{tension}: its pressure is below zero, like a stretched rope. This is the
\emph{cohesion--tension} explanation, the one accepted by plant physiologists. Its limit is also physical:
water under tension can suddenly form a vapour bubble (cavitation) that breaks the thread, and this risk,
growing with height, is thought to be one reason no tree grows much beyond about \SI{120}{m}.
\end{remarkbox}

\section{Summary}\label{sec:fluids-summary}
A fluid at rest can only push, and only perpendicular to a surface, so its forces are described by one
scalar, the pressure; a fluid in motion obeys conservation of mass (continuity) and of energy (Bernoulli);
viscosity and surface tension are the two corrections that real liquids demand.

\begin{itemize}
\item Pressure $P=F_\perp/A$ is a scalar, the same in every direction at a point. In a liquid at rest
      $P=P_0+\rho gh$, depending only on depth; $\rho gh$ is the gauge pressure. Air pressure falls
      exponentially with height, not linearly.
\item Pascal's law: a pressure change in an enclosed incompressible fluid is transmitted undiminished. A
      hydraulic lift gives $F_2=F_1A_2/A_1$ and conserves work, $F_1d_1=F_2d_2$.
\item Archimedes: the buoyant force equals the weight of the fluid displaced, for any shape and at any depth.
      A floating body displaces its own weight; its submerged fraction is $\rho_{\text{body}}/\rho_{\text{fluid}}$.
\item Continuity for an incompressible fluid: $A_1v_1=A_2v_2$.
\item Bernoulli, along a streamline of an ideal fluid: $P+\tfrac12\rho v^2+\rho gh=\text{constant}$, the
      work--energy theorem per unit volume. Applications: Torricelli's $v=\sqrt{2gh}$, the Venturi meter,
      lift (with care). It fails where viscosity, turbulence or compressibility matter.
\item Viscosity: $F=\eta A\,dv/dy$; $\eta$ in \si{Pa.s}. Flow in a pipe is laminar for
      $Re=\rho vD/\eta$ below about $2000$. Stokes' law $F=6\pi\eta rv$ (small, slow spheres, $Re\ll1$)
      gives the terminal velocity $v_t=2r^2(\rho_s-\rho)g/(9\eta)$, approached as
      $v_t(1-e^{-t/\tau})$.
\item Surface tension $S$ (\si{N/m}) is a force per unit length and an energy per unit area. A free drop
      is a sphere; excess pressure is $2S/R$ in a drop and $4S/R$ in a soap bubble.
\item Angle of contact: measured through the liquid to the surface's tangent; below $90^\circ$ the liquid
      wets and rises in a capillary, above $90^\circ$ it is depressed: $h=2S\cos\theta/(\rho gr)$.
\end{itemize}

\section{Problem Types}\label{sec:fluids-types}
Ten types cover the fluid questions examiners set. Most traps are of two kinds: forgetting that areas go as
the \emph{square} of a radius, and using a formula (Stokes, the capillary rise, the excess pressure) outside
the situation it was derived for.

\archq{Pressure at depth}
{``Find the gauge or absolute pressure at a given depth,'' or ``the force on a dam or a submarine
window.''}
{Gauge pressure $=\rho gh$; absolute pressure $=P_0+\rho gh$. Force on a small flat surface $=$ pressure
$\times$ area; for a tall vertical wall, integrate (the average gauge pressure is $\tfrac12\rho gH$).}
{Pressure depends only on the \emph{vertical depth}, never on the container's shape or the amount of fluid.
``Gauge'' excludes the atmosphere; ``absolute'' includes it.}
{The gauge pressure at a depth of \SI{10}{m} in a freshwater lake is about ($\rho=\SI{1000}{kg/m^3}$,
$g=\SI{9.8}{m/s^2}$):}
{\SI{9.8}{kPa}}
{\SI{98}{kPa}}
{\SI{980}{kPa}}
{\SI{199}{kPa}\archans{B}{$\rho gh=1000\times9.8\times10=\SI{98000}{Pa}=\SI{98}{kPa}$, close to one
atmosphere. Option (D) is the \emph{absolute} pressure, $98+101$~kPa; (A) and (C) slip a power of ten in
converting pascals to kilopascals.}}

\archq{The hydraulic lift}
{``A force on a small piston lifts a load on a large piston: find the load, or the force needed.''}
{The pressure is common: $F_1/A_1=F_2/A_2$, so $F_2=F_1A_2/A_1$. Areas go as the square of the radius (or
diameter).}
{Using the ratio of radii instead of the ratio of areas. And work is conserved: the large piston moves
less.}
{In a hydraulic lift the piston radii are in the ratio $1:10$. A force of \SI{50}{N} on the small piston can
support a load of:}
{\SI{500}{N}}
{\SI{50}{N}}
{\SI{5000}{N}}
{\SI{0.5}{N}\archans{C}{The force scales as the area, that is, as the radius squared: $10^2=100$, so
$F_2=50\times100=\SI{5000}{N}$. Option (A) uses the radius ratio; (B) confuses equal pressures with equal
forces; (D) inverts the ratio.}}

\archq{Floating and the submerged fraction}
{``What fraction of a floating body is submerged (or above the surface)?'' or ``what is its density?''}
{A floating body displaces its own weight: $V_{\text{sub}}/V=\rho_{\text{body}}/\rho_{\text{fluid}}$. The
fraction above is $1$ minus this.}
{Answering with the submerged fraction when the fraction above the surface is asked for, or dividing by the
wrong volume.}
{A block of ice ($\rho=\SI{917}{kg/m^3}$) floats in water ($\rho=\SI{1000}{kg/m^3}$). The fraction of its
volume above the surface is about:}
{$8\%$}
{$9\%$}
{$50\%$}
{$92\%$\archans{A}{Submerged fraction $=917/1000=0.917$, so the fraction above is $1-0.917\approx8\%$.
Option (D) is the submerged fraction; (B) divides the excess $1000-917$ by $917$ instead of $1000$; (C)
is a guess.}}

\archq{Apparent weight and relative density}
{``A body weighs $W_1$ in air and $W_2$ when fully immersed in water: find its density (or the liquid's).''}
{The loss of weight $W_1-W_2$ is the buoyant force, $\rho_wVg$. So
$\rho_{\text{body}}/\rho_w=W_1/(W_1-W_2)$.}
{Dividing by the weight in water, $W_2$, instead of by the \emph{loss} of weight.}
{A body weighs \SI{60}{N} in air and \SI{40}{N} when fully immersed in water. Its density is
($\rho_w=\SI{1000}{kg/m^3}$):}
{\SI{1500}{kg/m^3}}
{\SI{2000}{kg/m^3}}
{\SI{670}{kg/m^3}}
{\SI{3000}{kg/m^3}\archans{D}{Loss of weight $=\SI{20}{N}=\rho_wVg$; the density relative to water is
$60/20=3$, so $\rho=\SI{3000}{kg/m^3}$. Option (A) takes $60/40$; (B) takes $40/20$; (C) inverts
$40/60$.}}

\archq{Continuity: the narrowing pipe}
{``A pipe changes cross-section: find the new speed,'' often as the first step of a Bernoulli problem.}
{$A_1v_1=A_2v_2$: speed is inversely proportional to area, and area goes as diameter squared.}
{Halving the diameter quarters the area, so the speed rises fourfold, not twofold.}
{Water flows at \SI{2}{m/s} in a pipe. Where the pipe narrows to half its diameter, the speed becomes:}
{\SI{4}{m/s}}
{\SI{8}{m/s}}
{\SI{1}{m/s}}
{\SI{0.5}{m/s}\archans{B}{The area falls to a quarter, so the speed rises by $4$: $2\times4=\SI{8}{m/s}$.
Option (A) uses the diameter ratio; (C) and (D) have the fluid slowing down in the narrow part, against
continuity.}}

\archq{Torricelli's law}
{``Find the speed of water leaving a hole in a tank,'' or ``where the jet lands.''}
{Bernoulli from the free surface to the jet: both at $P_0$, $v_1\approx0$, so $v=\sqrt{2gh}$, with $h$ the
depth of the \emph{hole} below the surface. The jet is then a projectile.}
{Counting $\rho gh$ twice (pressure at the hole \emph{and} height of the surface), which gives
$\sqrt{4gh}$; or using the total depth of water instead of the depth of the hole.}
{Water flows from a small hole \SI{5.0}{m} below the surface of a large open tank. The speed of the jet is
about ($g=\SI{9.8}{m/s^2}$):}
{\SI{7.0}{m/s}}
{\SI{98}{m/s}}
{\SI{9.9}{m/s}}
{\SI{14}{m/s}\archans{C}{$v=\sqrt{2gh}=\sqrt{2\times9.8\times5.0}=\sqrt{98}\approx\SI{9.9}{m/s}$. Option
(A) drops the $2$, $\sqrt{gh}$; (B) forgets the square root; (D) counts $\rho gh$ twice,
$\sqrt{4gh}$.}}

\archq{Bernoulli in a horizontal pipe}
{``The speed changes between two points at the same height: find the pressure difference,'' or ``find the
flow rate from a measured pressure difference'' (Venturi meter).}
{At one height $P_1-P_2=\tfrac12\rho(v_2^2-v_1^2)$; get the speeds from continuity first.}
{Using $\tfrac12\rho(v_2-v_1)^2$ instead of $\tfrac12\rho(v_2^2-v_1^2)$, or dropping the $\tfrac12$. Fast flow
has the \emph{lower} pressure.}
{Water flows through a horizontal pipe at \SI{2.0}{m/s}; in a narrow section its speed is \SI{6.0}{m/s}. The
pressure in the narrow section is lower by:}
{\SI{16}{kPa}}
{\SI{32}{kPa}}
{\SI{8}{kPa}}
{\SI{18}{kPa}\archans{A}{$\tfrac12(1000)(6.0^2-2.0^2)=500\times32=\SI{16000}{Pa}$. Option (B) drops the
$\tfrac12$; (C) squares the difference of the speeds, $\tfrac12\rho(4.0)^2$; (D) forgets the slow flow's
kinetic term, $\tfrac12\rho v_2^2$.}}

\archq{Terminal velocity and Stokes' law}
{``Find the terminal velocity of a small sphere in a viscous fluid,'' or ``how it changes when drops
merge or the radius changes.''}
{$v_t=2r^2(\rho_s-\rho)g/(9\eta)$, so $v_t\propto r^2$. When drops merge, the volume is conserved: $n$
drops of radius $r$ make one of radius $n^{1/3}r$.}
{Scaling with $r$ instead of $r^2$, or adding the volumes as if the radius added.}
{Eight identical small raindrops, each falling at terminal speed $v$ in the Stokes regime, merge into one
drop. Its terminal speed is:}
{$4v$}
{$2v$}
{$8v$}
{$v$\archans{A}{Volume conserved: $\tfrac43\pi R^3=8\times\tfrac43\pi r^3$, so $R=2r$; and $v_t\propto r^2$
gives $4v$. Option (B) scales with $r$; (C) scales with the mass; (D) forgets that a larger drop falls
faster.}}

\archq{Excess pressure in drops and bubbles}
{``Find the pressure inside a drop or a bubble,'' or ``which way does air flow between two bubbles.''}
{Drop or air bubble inside a liquid (one surface): $2S/R$. Soap bubble in air (two surfaces): $4S/R$.
Convert the radius to metres.}
{Using the diameter as the radius, or the one-surface formula for a soap bubble.}
{A soap bubble in air has a diameter of \SI{2.0}{cm}. If $S=\SI{0.030}{N/m}$, the excess pressure inside
it is:}
{\SI{3}{Pa}}
{\SI{6}{Pa}}
{\SI{0.12}{Pa}}
{\SI{12}{Pa}\archans{D}{$R=\SI{0.010}{m}$ and two surfaces: $4S/R=4(0.030)/0.010=\SI{12}{Pa}$. Option
(B) uses $2S/R$ (one surface) or puts the diameter for the radius in $4S/R$; (A) makes both slips at once;
(C) leaves the radius in centimetres.}}

\archq{Capillary rise}
{``Find the height a liquid rises (or falls) in a capillary,'' or ``how the height changes with the radius or
with tilting.''}
{$h=2S\cos\theta/(\rho gr)$, so $h\propto1/r$. For a tilted tube the \emph{vertical} height stays the same,
so the length of liquid in the tube grows.}
{Scaling with $1/r^2$, or expecting the liquid to overflow a short tube.}
{Water rises \SI{6.0}{cm} in a clean glass capillary. In a tube of half the radius it rises:}
{\SI{3.0}{cm}}
{\SI{12}{cm}}
{\SI{24}{cm}}
{\SI{6.0}{cm}\archans{B}{$h\propto1/r$, so halving $r$ doubles $h$ to \SI{12}{cm}. Option (A) takes
$h\propto r$; (C) takes $h\propto1/r^2$, as if the rise went with the area; (D) assumes the rise does not
depend on the radius.}}

\section{Practice Problems}\label{sec:fluids-practice}

The problems below are of standard examination difficulty. Take $g=\SI{9.8}{m/s^2}$,
$P_0=\SI{1.013e5}{Pa}$ and the density of water \SI{1000}{kg/m^3} unless told otherwise. Answers, with brief
working, follow the list.

\begin{enumerate}[itemsep=3pt]
\item Find the gauge and the absolute pressure at the bottom of a swimming pool \SI{2.0}{m} deep.
\item A U-tube contains water. Oil of density \SI{800}{kg/m^3} is poured into one arm until it stands
      \SI{12}{cm} above the oil--water boundary. How high does the water in the other arm stand above that
      boundary?
\item A wooden block floats in water with $60\%$ of its volume submerged. Find its density, and the fraction
      submerged when it floats in oil of density \SI{800}{kg/m^3}.
\item The pistons of a hydraulic lift have areas \SI{5.0}{cm^2} and \SI{250}{cm^2}. What force on the small
      piston supports a load of \SI{1000}{kg}?
\item Water flows through a horizontal pipe at \SI{1.5}{m/s} with pressure \SI{2.0e5}{Pa}. Find the pressure
      where the pipe's cross-section is reduced to a third.
\item An open tank has a hole of area \SI{2.0}{cm^2} at \SI{1.25}{m} below the water surface. Find the speed
      of the jet and the volume of water leaving each second (ignore the contraction of the jet).
\item A steel ball of radius \SI{1.0}{mm} ($\rho=\SI{7800}{kg/m^3}$) falls through glycerine
      ($\rho=\SI{1260}{kg/m^3}$, $\eta=\SI{1.5}{Pa.s}$). Find its terminal velocity, and check that Stokes'
      law applies.
\item How much work is needed to blow a soap bubble of radius \SI{3.0}{cm}, if $S=\SI{0.030}{N/m}$?
\item A drop of water of radius \SI{1.0}{mm} is sprayed into $1000$ identical droplets. Find the energy
      needed ($S=\SI{0.072}{N/m}$).
\item A clean glass capillary of inner radius \SI{0.25}{mm} stands in water ($S=\SI{0.072}{N/m}$,
      $\theta=0$). How high does the water rise? What happens if only \SI{4.0}{cm} of the tube projects
      above the water?
\end{enumerate}

\subsection*{Answers}
{\small
\begin{enumerate}[itemsep=2pt]
\item Gauge: $\rho gh=1000(9.8)(2.0)=\SI{1.96e4}{Pa}$. Absolute: $1.013\times10^5+0.196\times10^5
\approx\SI{1.21e5}{Pa}$.
\item The pressures at the level of the boundary are equal in the two arms (same level, same connected
liquid below it), so $\rho_{\text{oil}}(12)=\rho_w h$ and $h=0.8\times12=\SI{9.6}{cm}$.
\item $\rho=0.60\times1000=\SI{600}{kg/m^3}$. In oil the submerged fraction is $600/800=0.75$, that is,
$75\%$.
\item $F_1=mg\,A_1/A_2=1000(9.8)(5.0/250)=\SI{196}{N}$.
\item Continuity: $v_2=3\times1.5=\SI{4.5}{m/s}$. Bernoulli: $P_2=2.0\times10^5-\tfrac12(1000)(4.5^2-1.5^2)
=2.0\times10^5-9000=\SI{1.91e5}{Pa}$.
\item $v=\sqrt{2(9.8)(1.25)}=\sqrt{24.5}\approx\SI{4.95}{m/s}$; flow rate
$Av=(2.0\times10^{-4})(4.95)\approx\SI{9.9e-4}{m^3/s}$, about \SI{1}{\litre} each second.
\item $v_t=2r^2(\rho_s-\rho)g/(9\eta)=2(10^{-3})^2(7800-1260)(9.8)/(9\times1.5)$, which is about
\SI{9.5e-3}{m/s}, just under \SI{1}{cm/s}. Check: $Re=\rho vD/\eta=1260(0.0095)(0.002)/1.5$, which is
about $0.016$, far below $1$, so Stokes' law holds.
\item Two surfaces: $W=S\times2(4\pi R^2)=0.030\times8\pi(0.030)^2\approx\SI{6.8e-4}{J}$.
\item Volume conserved: $1000r^3=R^3$, so $r=R/10=\SI{0.10}{mm}$. New area minus old:
$4\pi(1000r^2-R^2)=4\pi(10^{-5}-10^{-6})\approx\SI{1.13e-4}{m^2}$. Energy $=S\Delta A\approx\SI{8.1e-6}{J}$.
\item $h=\dfrac{2(0.072)}{1000(9.8)(2.5\times10^{-4})}\approx\SI{5.9}{cm}$. With only \SI{4.0}{cm} of tube,
the water rises to the top and stops; the meniscus flattens (its radius of curvature grows to about
$5.9/4.0\approx1.5$ times the tube's radius) so that its smaller excess pressure supports the
\SI{4.0}{cm} column. It does not overflow.
\end{enumerate}
}

\section{Challenge Problems}\label{sec:fluids-challenge}

The problems below are harder than examination questions, chiefly because each needs a principle of the
chapter applied in a setting where no formula has been prepared for it. Attempt each fully before reading its
solution.

\begin{figure}[htbp]
\centering
\tikzfig{fluids_bucket}
\figcaption{fig:fluids-bucket}{Water in a spinning bucket, drawn to the numbers of the challenge. In the frame
turning with the bucket, each parcel feels $g$ down and the pseudo-force per unit mass $\omega^2r$ outwards;
the free surface sets itself perpendicular to their resultant (dashed), which makes it a paraboloid. The
level at rest lies halfway between the centre and the rim of the curved surface.}
\end{figure}

\begin{challengebox}{The rotating bucket (builds on pressure at depth)}
A cylindrical bucket of water of radius $R=\SI{10}{cm}$ spins steadily about its vertical axis at
$\omega=\SI{10}{rad/s}$, the water turning with it. By how much does the water at the rim stand above the
water at the axis?

\challengeans{Work in the frame rotating with the bucket (\chref{ch:circdyn},
Section~\ref{sec:noninertial}), in which the water is at rest. Each parcel at distance $r$ from the axis feels
gravity $g$ down and a centrifugal pseudo-force $\omega^2r$ per unit mass outwards. A free surface at rest
must be perpendicular to the net effective gravity (otherwise the surface liquid would slide along it), so
its slope is
\[
\frac{dz}{dr}=\frac{\omega^2r}{g}\quad\Longrightarrow\quad z(r)=\frac{\omega^2r^2}{2g},
\]
a paraboloid (Figure~\ref{fig:fluids-bucket}). The rise at the rim is
\[
\Delta z=\frac{\omega^2R^2}{2g}=\frac{(10)^2(0.10)^2}{2(9.8)}\approx\SI{51}{mm}.
\]
In the ground frame the same result follows without pseudo-forces: a ring of water at radius $r$ moves in a
circle, so the horizontal pressure gradient must supply its centripetal acceleration,
$dP/dr=\rho\omega^2r$; along the surface the pressure is $P_0$, so the depth below the surface must grow
outwards to match. \ans{\SI{51}{mm}}}
\end{challengebox}

\begin{figure}[htbp]
\centering
\tikzfig{fluids_acceltank}
\figcaption{fig:fluids-acceltank}{Water in a tank accelerating to the right, drawn to the numbers of the
challenge. In the tank's frame the effective gravity is $g$ down plus the pseudo-force per unit mass $a$
backwards; the surface settles perpendicular to it, tilted at $\tan\theta=a/g$ and deeper at the rear.}
\end{figure}

\begin{challengebox}{The accelerating tank (builds on pressure at depth)}
An open tank of water is carried on a truck accelerating horizontally at $a=\SI{2.0}{m/s^2}$. At what angle to
the horizontal does the water surface settle, and which end is deeper?

\challengeans{In the truck's frame the water is at rest and feels gravity $g$ down plus a pseudo-force $a$ per
unit mass backwards (\chref{ch:circdyn}). The free surface sets itself perpendicular to their resultant
(Figure~\ref{fig:fluids-acceltank}), tilting by
\[
\tan\theta=\frac{a}{g}=\frac{2.0}{9.8}\quad\Longrightarrow\quad\theta\approx11.5^\circ,
\]
with the water deeper at the \emph{rear}. In the ground frame: the water must be accelerated forwards, which
needs a higher pressure behind each parcel than in front, and the extra depth at the rear supplies it. The
same $\tan\theta=a/g$ gives the angle of a plumb line hanging in an accelerating train. \ans{$11.5^\circ$,
deeper at the rear}}
\end{challengebox}

\begin{challengebox}{Capillarity between plates (builds on capillary rise)}
Two parallel, clean glass plates a distance $d=\SI{0.20}{mm}$ apart are dipped vertically into water
($S=\SI{0.072}{N/m}$, $\theta\approx0$). How high does the water climb between them? Compare with a
capillary tube whose \emph{diameter} equals the gap.

\challengeans{Take a width $w$ of the plates (along the water surface). The water touches each plate along a
line of length $w$, so two lines pull up with $2Sw\cos\theta$. They hold up a slab of width $w$, thickness
$d$ and height $h$, weighing $\rho gwdh$. So
\[
2Sw\cos\theta=\rho gwdh\quad\Longrightarrow\quad h=\frac{2S\cos\theta}{\rho gd}
=\frac{2(0.072)}{(1000)(9.8)(0.20\times10^{-3})}\approx\SI{73}{mm}.
\]
A tube of diameter $d$ has radius $r=d/2$, and $h=2S/(\rho gr)=4S/(\rho gd)\approx\SI{147}{mm}$, twice as
high. The reason is geometric: per unit length of the pulling edge, the plates must hold up a slab of
thickness $d/2$ (area $wd$ shared by two edges of length $w$), while the tube holds up only
$\pi r^2/(2\pi r)=r/2=d/4$. (The plates rise as high as a tube whose \emph{radius} equals the gap.)
\ans{\SI{73}{mm}, half the tube's}}
\end{challengebox}

\begin{figure}[htbp]
\centering
\tikzfig{fluids_bubbles}
\figcaption{fig:fluids-bubbles}{Two soap bubbles, drawn to scale for $r_1=\SI{2.0}{cm}$ and $r_2=\SI{3.0}{cm}$.
The smaller bubble has the higher pressure, so the shared film bulges into the larger one, with radius of
curvature $R=\SI{6.0}{cm}$. The three films meet at equal angles of $120^\circ$, which fixes how far apart the
centres are.}
\end{figure}

\begin{challengebox}{Two soap bubbles in contact (builds on excess pressure)}
Two soap bubbles of radii $r_1=\SI{2.0}{cm}$ and $r_2=\SI{3.0}{cm}$ meet and share a common film. What is the
radius of curvature of that film, and which way does it bulge?

\challengeans{Each bubble's interior is at an excess pressure $4S/r$ above the air outside. The shared film
separates the two interiors, so it carries the difference,
$4S\big(\tfrac1{r_1}-\tfrac1{r_2}\big)$. A film of radius of curvature $R$ supports $4S/R$, so
\[
\frac1R=\frac1{r_1}-\frac1{r_2}\quad\Longrightarrow\quad R=\frac{r_1r_2}{r_2-r_1}=\frac{(2.0)(3.0)}{3.0-2.0}
=\SI{6.0}{cm}.
\]
The smaller bubble has the higher pressure, so the film bulges into the larger bubble
(Figure~\ref{fig:fluids-bubbles}). For equal bubbles $R\to\infty$: the shared film is flat.
\ans{\SI{6.0}{cm}, bulging into the larger bubble}}
\end{challengebox}

\begin{challengebox}{The raindrop that cannot obey Stokes (builds on terminal velocity)}
Estimate the terminal velocity of a raindrop of radius $r=\SI{1.0}{mm}$ from Stokes' law
($\eta_{\text{air}}=\SI{1.8e-5}{Pa.s}$, $\rho_{\text{air}}=\SI{1.2}{kg/m^3}$). Then compute the Reynolds number
$Re=\rho_{\text{air}}v(2r)/\eta$ and explain why the answer is wrong. Finally, estimate the true terminal
velocity using a drag force $\tfrac12C\rho_{\text{air}}Av^2$, with $A=\pi r^2$ and a drag coefficient
$C\approx0.5$.

\challengeans{Stokes' law gives
\[
v_t=\frac{2r^2(\rho_w-\rho_{\text{air}})g}{9\eta}\approx\frac{2(10^{-3})^2(1000)(9.8)}{9(1.8\times10^{-5})}
\approx\SI{121}{m/s},
\]
which is absurd: raindrops of this size fall at about \SI{6.5}{m/s}. The Reynolds number shows why:
\[
Re=\frac{(1.2)(121)(2\times10^{-3})}{1.8\times10^{-5}}\approx1.6\times10^4,
\]
far beyond the $Re\ll1$ that Stokes' law requires. At such speeds the drop leaves a turbulent wake, and the
drag is dominated by the momentum given to the air it pushes aside, which grows as $v^2$. Balancing that drag
against the weight, $\tfrac12C\rho_{\text{air}}\pi r^2v^2=\tfrac43\pi r^3\rho_wg$, gives
\[
v=\sqrt{\frac{8r\rho_wg}{3C\rho_{\text{air}}}}=\sqrt{\frac{8(10^{-3})(1000)(9.8)}{3(0.5)(1.2)}}
\approx\SI{6.6}{m/s},
\]
close to the measured value. The lesson: check $Re$ before trusting Stokes' law. It holds for fog droplets
and fine dust, not for raindrops. \ans{Stokes \SI{121}{m/s} (wrong); about \SI{6.6}{m/s}}}
\end{challengebox}

\section*{Looking Back on Part A}
\addcontentsline{toc}{section}{Looking Back on Part A}

\textbf{Part A applied the mechanics of Volume~I to bulk matter.} Forces and equilibrium
(\chref{ch:newton}) and work and energy (\chref{ch:energy}) were all it needed; what was new was the way a
material responds.
\begin{itemize}
\item \textbf{Solids (\chref{ch:solids}).} A solid resists deformation. The response is described by
\emph{stress} (force per unit area) and \emph{strain} (fractional deformation), tied together by a modulus
for each kind of deformation: Young's modulus $Y$ for stretching, the bulk modulus $B$ for squeezing, and the
shear modulus $G$ for shearing. Below the elastic limit a solid behaves as a spring and stores
$\tfrac12\times\text{stress}\times\text{strain}$ of energy per unit volume; beyond it, it yields and
eventually breaks.
\item \textbf{Fluids at rest (\chref{ch:fluids}).} A fluid has no shear modulus, so at rest it can only push
perpendicular to a surface, and its forces reduce to one scalar, pressure. Pressure grows with depth as
$P=P_0+\rho gh$, is transmitted undiminished through an enclosed liquid (Pascal, hence hydraulics), and
buoys a body with the weight of the fluid it displaces (Archimedes, hence floating).
\item \textbf{Fluids in motion (\chref{ch:fluids}).} Continuity ($Av=\text{constant}$) is conservation of
mass; Bernoulli's equation ($P+\tfrac12\rho v^2+\rho gh=\text{constant}$) is the work--energy theorem per unit
volume. Viscosity, a shear stress proportional to the \emph{rate} of strain, and surface tension, an energy
per unit area of surface, are the corrections real fluids demand.
\end{itemize}
The two chapters share one language. Stress and pressure are both force per unit area, measured in pascals;
the bulk modulus of a solid and the near-incompressibility of a liquid are the same property; and the
energy stored in a stretched wire, in a raised column of liquid and in a liquid surface are all energies
per unit volume or per unit area, traded against work.

\textbf{The thread into Part B:} throughout Part A the temperature was fixed. Yet the moduli, the viscosity
and the surface tension all change when a material is heated, and a gas's pressure depends on its
temperature as much as on its volume. \chref{ch:thermal} lets the temperature change, and finds that it
governs expansion, changes of state and the flow of heat itself.
